% Chapter 4: 실험 결과
\chapter{실험 결과}
\label{ch:results}

본 장에서는 제~\ref{ch:mutbench}장에서 제시한 MutBench 프레임워크의 실험 결과를 2단계 구조로 보고한다.

\textbf{1단계 (깊이 우선, 제~\ref{sec:variant_comparison}--\ref{sec:phylo_simulation}절):}
SARS-CoV-2 단일 병원체에서 5가지 MutClust 변이체와 2가지 기준선 방법의 성능을 정밀 비교한다.
합성 벤치마크, DMS 검증, 민감도 분석, 시간적 일반화, 다중 ground truth 평가, 교차 병원체 예비 검증(H3N2)을 통해, ``탐지 전략의 선택이 파라미터 조정보다 성능에 결정적이며, 최적 방법이 데이터 분포 특성에 의존한다''는 교훈을 도출한다.

\textbf{2단계 (너비 우선, 제~\ref{sec:9pathogen_results}절):}
1단계의 교훈을 일반화하기 위해, 9종 병원체 $\times$ 9가지 스코어링 $\times$ 41개 탐지 방법으로 확장한 대규모 벤치마크를 수행한다.
1단계의 MutClust 변이체들은 단일 알고리즘의 구조적 변형이므로, 2단계에서는 원본(MutClust-Original)만을 대표로 포함하고 다양한 알고리즘 패밀리(Wavelet, KDE, CUSUM, Bayesian 등)의 독립 방법과 비교한다.

%% ============================================================
\section{합성 벤치마크: 변이체 비교}
\label{sec:variant_comparison}

합성 H-score를 사용하여 SARS-CoV-2 유사 게놈(29,903 위치)에서 7개 기능 영역(1,251개 고유 ground truth 뉴클레오티드 위치)을 대상으로 방법을 비교하였다.
제~\ref{ch:mutbench}장에서 기술한 바와 같이, 이 합성 벤치마크는 통제된 신호 복원 검정으로서 알려진 위치에 삽입된 신호의 복원 능력을 평가한다.
결과는 표~\ref{tab:method_comparison}에 제시하였다.

\begin{table}[htbp]
\centering
\caption{합성 SARS-CoV-2 데이터에 대한 방법 비교 (29,903 위치, 1,251개 ground truth 기능 뉴클레오티드 위치)}
\label{tab:method_comparison}
\begin{tabular}{lrrrrrr}
\toprule
\textbf{방법} & \textbf{정밀도} & \textbf{재현율} & \textbf{F1} & \textbf{농축비} & \textbf{탐지 수} & \textbf{클러스터} \\
\midrule
MutClust-Original & 0.991 & 0.504 & 0.668 & 23.68 & 637 & 31 \\
MutClust-Fixed & 0.988 & 0.506 & 0.669 & 23.60 & 641 & 30 \\
MutClust-dNdS$^\dagger$ & 0.988 & 0.506 & 0.669 & 23.60 & 641 & 30 \\
MutClust-Hybrid & 0.968 & 0.659 & 0.785 & 23.15 & 852 & 8 \\
HDBSCAN-Auto & 0.079 & 0.944 & 0.146 & 1.89 & 14,916 & 21 \\
\midrule
빈도 기준선 & 0.404 & 0.965 & 0.569 & 9.65 & 2,991 & 1,587 \\
슬라이딩 윈도우 & 0.157 & 0.974 & 0.270 & 3.75 & 7,773 & 214 \\
\bottomrule
\end{tabular}
\begin{minipage}{\textwidth}
\smallskip
\footnotesize $^\dagger$합성 데이터에서는 dN/dS 계산을 위한 참조 서열이 제공되지 않으므로, MutClust-Fixed와 동일한 결과를 산출한다.
\end{minipage}
\end{table}

결과는 탐지 전략 간의 명확한 정밀도--재현율 절충을 드러낸다.
CCM 기반 시드 탐지를 사용하는 MutClust 변이체(Original, Fixed, dNdS)는 거의 완벽한 정밀도($> 0.98$)를 달성하지만 중간 수준의 재현율(0.50--0.51)을 보이며, 이는 CCM 탐지가 주로 고신호 핫스팟 영역(RBD, 퓨린 절단)을 식별하면서 저신호 기능 부위(NTD 초과부위, HR1/HR2)를 놓치는 것을 반영한다.
버그 수정은 정밀도를 유지하면서 탐지 위치를 4개 확장하였다(641 vs. 637).

MutClust-Hybrid는 CCM 시드 탐지와 HDBSCAN 경계 결정을 결합하여 최고 F1(0.785)을 달성하였다.
정밀도(0.968)는 CCM 전용 변이체에 근접하면서 재현율(0.659)이 크게 향상되었으며, HDBSCAN의 더 넓은 경계 추정을 통해 NTD 초과부위의 일부를 포함한 추가 기능 영역을 포착하였다.

HDBSCAN-Auto는 최고 재현율(0.944)을 달성하였으나 낮은 정밀도(0.079)를 보이며, 14,916개 위치---게놈의 약 절반---를 21개 클러스터에서 탐지하였다.
CCM 시드 탐지가 후보 공간을 제한하지 않으면 HDBSCAN이 실제 기능 부위와 함께 상당한 배경 잡음을 포착함을 보여준다.
이는 도메인 특이적 사전 지식이 탐지 특이성을 실질적으로 향상시킴을 입증한다.

\textbf{외부 기준선 비교.}
빈도 기준선(H-score 역치 초과 위치 탐지)과 슬라이딩 윈도우(고정 크기 창의 평균 H-score 역치)는 모두 높은 재현율(0.965, 0.974)을 달성하지만 낮은 정밀도(0.404, 0.157)를 보인다.
이 기준선들의 F1(0.569, 0.270)은 모든 MutClust 변이체(0.668--0.785)보다 낮으며, 특히 MutClust-Hybrid 대비 F1 격차가 0.216--0.515에 달한다.
빈도 기준선은 적응적 클러스터 크기 결정 능력이 없어 2,991개 위치를 1,587개의 독립 ``클러스터''로 탐지하며, 슬라이딩 윈도우는 고정 창 크기의 제약으로 7,773개 위치를 과잉 탐지한다.
이 결과는 MutClust의 CCM 기반 적응형 밀도 추정이 단순 역치·윈도우 기반 접근법 대비 실질적인 성능 향상을 제공함을 확인한다.

Hotspot-score를 포함한 확장된 평가 지표는 표~\ref{tab:hotspot_scores}에 제시하였다.

\begin{table}[htbp]
\centering
\caption{Hotspot-score 분석 (100회 부분추출, 80\% 위치 유지)}
\label{tab:hotspot_scores}
\begin{tabular}{lrrrr}
\toprule
\textbf{방법} & \textbf{정밀도} & \textbf{재현율} & \textbf{안정성} & \textbf{Hotspot-score} \\
\midrule
MutClust-Original & 0.991 & 0.504 & 0.418 & 0.638 \\
MutClust-Fixed & 0.988 & 0.506 & 0.445 & 0.646 \\
MutClust-dNdS & 0.988 & 0.506 & 0.445 & 0.646 \\
\textbf{MutClust-Hybrid} & \textbf{0.968} & \textbf{0.659} & \textbf{0.706} & \textbf{0.778} \\
HDBSCAN-Auto & 0.079 & 0.944 & 0.788 & 0.604 \\
\bottomrule
\end{tabular}
\end{table}

Hotspot-score는 F1만으로는 포착할 수 없는 평가 차원을 드러내며, \textbf{MutClust-Hybrid를 최우수 방법}(hotspot-score 0.778)으로 식별한다.
MutClust-Hybrid의 우위는 세 차원 모두에서의 강한 성능에서 비롯된다: 거의 완벽한 정밀도(0.968), 향상된 재현율(0.659), 높은 안정성(0.706).
CCM 시드 탐지를 유지하면서 HDBSCAN을 경계 결정에 사용함으로써, 순수 CCM 탐지가 놓치는 기능 영역(예: 저신호 NTD 초과부위 위치)을 포착하면서 도메인 특이적 사전 지식의 특이성 이점을 보존한다.

HDBSCAN-Auto는 최고 개별 안정성(0.788)과 재현율(0.944)을 달성하지만, 낮은 정밀도(0.079)로 인해 MutClust 변이체 중 최저 hotspot-score(0.604)를 기록한다.
이는 복합 지표의 가치를 잘 보여준다: 재현율만으로 HDBSCAN-Auto를 평가하면 최우수 방법으로 순위가 매겨지고, 정밀도만으로 평가하면 완전히 기각될 것이다.
Hotspot-score는 하이브리드 접근법을 최적 균형으로 올바르게 식별하는 균형 잡힌 평가를 제공한다(그림~\ref{fig:mutbench_method_comparison}).

\textbf{2단계 파이프라인 검증.}
HDBSCAN-Auto의 높은 재현율과 MutClust-Hybrid의 높은 정밀도를 결합하기 위해, 2단계 탐지 파이프라인을 추가적으로 검증하였다.
HDBSCAN-Auto가 탐지한 14,916개 위치와 MutClust-Hybrid가 탐지한 852개 위치의 교집합(831개 위치, 6개 클러스터)은 정밀도 0.968을 유지하면서 안정성이 0.740으로 향상되어 최고 hotspot-score(0.784)를 달성하였으며, F1(0.772)은 MutClust-Hybrid 단독(0.785)에 근접하였다(상세 결과는 제~\ref{ch:discussion}장 표~\ref{tab:twostep_pipeline} 참조).
반면 합집합(14,964개 위치)은 재현율 0.970과 안정성 0.794에서 최고를 기록하였으나 정밀도가 0.081로 급락하였다.
이 결과는 두 방법의 상보성이 교집합 전략에서 가장 효과적으로 발현됨을 입증하며, 응용 맥락별 전략 선택의 근거를 제공한다(제~\ref{ch:discussion}장에서 상세 논의).

\begin{figure}[htbp]
\centering
\includegraphics[width=0.9\textwidth]{method_comparison.png}
\caption{합성 SARS-CoV-2 데이터에 대한 방법 비교. MutClust-Hybrid가 정밀도와 재현율의 최적 균형(F1 = 0.785)을 달성하며, CCM 전용 변이체는 높은 정밀도이지만 중간 재현율을, HDBSCAN-Auto는 높은 재현율이지만 매우 낮은 정밀도를 보인다.}
\label{fig:mutbench_method_comparison}
\end{figure}

%% ============================================================
\section{통계적 검증}
\label{sec:statistical_validation}

벤치마크 결과의 통계적 견고성을 평가하기 위해 부트스트랩 신뢰 구간 추정, 순열 검정(무작위 레이블 섞기 + 순환 순열), 쌍별 비교 검정, Fisher's exact test를 수행하였다.
본 절에서 보고하는 모든 부트스트랩 신뢰 구간은 소표본($n = 7$)에서의 편향과 비대칭을 보정하기 위해 bias-corrected and accelerated (BCa) 방법~\cite{efron1987better}을 사용하여 추정하였다.

\textbf{순열 검정.}
무작위 레이블 섞기(1,000회)에서 모든 방법의 정밀도·재현율이 $p < 0.001$로 유의하였다.
순환 순열 검정(200회)에서도 모든 5개 방법의 관찰 F1이 귀무 분포를 크게 초과하였다:
MutClust-Hybrid(F1 = 0.785, 귀무 평균 = 0.023, $z$ = 9.97, $p < 0.005$),
MutClust-Original/Fixed(F1 = 0.669, 귀무 평균 = 0.018, $z$ = 9.25, $p < 0.005$),
HDBSCAN-Auto(F1 = 0.146, 귀무 평균 = 0.076, $z$ = 11.31, $p < 0.005$).
H-score의 공간 자기상관을 보존하는 순환 순열에서도 모든 방법이 유의함은, 탐지된 핫스팟과 기능 영역 간의 공간적 중첩이 H-score 클러스터 구조의 부산물이 아닌 진정한 생물학적 신호임을 확인한다.

\textbf{Hotspot-score 귀무 분포.}
복합 지표인 hotspot-score 자체의 통계적 유의성을 검증하기 위해, 동일한 순환 순열(200회)을 통해 hotspot-score의 귀무 분포를 추정하였다.
모든 5개 방법의 관찰 hotspot-score가 귀무 분포를 유의하게 초과하였다($p < 0.005$):
MutClust-Hybrid(HS = 0.778, 귀무 평균 = 0.251),
MutClust-Fixed(HS = 0.646, 귀무 평균 = 0.162),
MutClust-Original(HS = 0.638, 귀무 평균 = 0.153),
HDBSCAN-Auto(HS = 0.604, 귀무 평균 = 0.441).
F1과 hotspot-score 모두 동일한 순환 순열에서 산출되므로 방법 간 순위 패턴이 일관되게 나타난다.
이 결과는 hotspot-score의 세 구성 요소(재현율, 정밀도, 안정성) 중 공간적 중첩에 의존하는 재현율과 정밀도가 무작위 배치에서는 달성될 수 없는 수준임을 확인하며, 복합 지표의 관찰값이 우연에 기인하지 않음을 입증한다.

\textbf{Fisher's exact test.}
다중 ground truth 평가(제~\ref{sec:multi_gt}절)에서 보고된 농축비의 통계적 유의성을 Fisher's exact test로 검증하였다.
128개 비교에 대해 Benjamini-Hochberg FDR 보정($q < 0.05$)을 적용한 결과, 3개 비교가 유의하였다:
TC-freq-v2의 기능 영역 농축(1.16배, $p_{\text{adj}}$ = 0.002),
TC-freq-v3의 기능 영역 농축(1.59배, $p_{\text{adj}}$ = 0.023),
CH-score의 수렴 진화 농축(1.70배, $p_{\text{adj}}$ = 0.023).
추가로 2개 비교가 $q < 0.10$에서 유의하였다(CH-score 수렴 진화 top10pct, E-score 수렴 진화 top20pct).
TC-freq-v3의 수렴 진화 엄격 부분집합 농축(3.43배, $p$ = 0.013, $p_{\text{adj}}$ = 0.161)과 포화 H-score의 수렴 진화 농축(1.85배, $p$ = 0.048, $p_{\text{adj}}$ = 0.269)은 BH FDR 보정 후 유의 수준에 미치지 못하나, 수렴 진화 위치에서의 농축 경향은 복수의 접근법(CH-score, E-score, TC-freq-v3)에서 일관되게 관찰된다.

\textbf{BCa 부트스트랩 신뢰 구간.}
$n = 7$ 기능 영역이라는 작은 표본에서 백분위 신뢰 구간은 편향과 비대칭을 적절히 반영하지 못할 수 있다.
이를 보정하기 위해 bias-corrected and accelerated (BCa) 부트스트랩 신뢰 구간을 적용하였다.
BCa 방법은 편향 보정 인자 $z_0$(부트스트랩 분포의 중위수 편향)와 가속 인자 $a$(잭나이프 기반 왜도 보정)를 사용하여 신뢰 구간의 백분위를 조정한다~\cite{efron1987better}.
7개 기능 영역을 복원추출(1,000회)하여 영역별 재현율·정밀도·안정성을 계산한 후 hotspot-score의 BCa 95\% 신뢰 구간을 표~\ref{tab:bootstrap_ci}에 제시하였다.

\begin{table}[htbp]
\centering
\caption{Hotspot-score의 영역 수준 BCa 부트스트랩 95\% 신뢰 구간 (7개 기능 영역의 1,000회 리샘플링, bias-corrected and accelerated)}
\label{tab:bootstrap_ci}
\begin{tabular}{lrrrcc}
\toprule
\textbf{방법} & \textbf{HS$_{\text{region}}$} & \textbf{BCa 하한} & \textbf{BCa 상한} & \textbf{$z_0$} & \textbf{$a$} \\
\midrule
\textbf{MutClust-Hybrid} & \textbf{0.476} & \textbf{0.323} & \textbf{0.635} & $-$0.090 & 0.010 \\
HDBSCAN-Auto & 0.468 & 0.339 & 0.567 & $-$0.003 & $-$0.025 \\
MutClust-Fixed & 0.311 & 0.196 & 0.534 & 0.030 & 0.070 \\
MutClust-Original & 0.302 & 0.187 & 0.505 & $-$0.015 & 0.070 \\
\bottomrule
\end{tabular}
\begin{minipage}{\textwidth}
\smallskip
\footnotesize HS$_{\text{region}}$: 7개 영역별 hotspot-score의 평균. 전체 집계 hotspot-score(표~\ref{tab:hotspot_scores})와 달리 영역 수준 평균은 고신호·저신호 영역에 동일 가중치를 부여하므로 절대값이 낮아진다.
\end{minipage}
\end{table}

영역 수준 BCa 부트스트랩은 7개 기능 영역을 복원추출하여 ground truth 정의의 불확실성---즉, ``어떤 기능 영역을 포함하느냐에 따라 평가 결과가 어떻게 달라지는가''---를 정량화한다.
CCM 기반 방법(MutClust-Original, Fixed)의 양의 가속 인자($a = 0.070$)는 영역별 점수 분포가 오른쪽으로 치우쳐져 있음을 나타내며, 이는 RBD(고신호 영역) 포함 여부에 따른 점수 변동을 반영한다.
반면 HDBSCAN-Auto의 좁은 구간(0.339--0.567)과 거의 0인 가속 인자($a = -0.025$)는 광범위한 탐지 범위가 영역 구성 변화에 덜 민감하고 대칭적 분포를 가짐을 보여준다.
MutClust-Hybrid의 BCa 구간(0.323--0.635)은 백분위 구간(0.332--0.651)에 비해 하한이 약간 낮고 상한이 약간 낮아, 편향 보정의 효과가 관찰된다.
모든 방법에서 점추정치가 95\% BCa CI 내에 위치하여, 부트스트랩 설계의 통계적 유효성이 확보되었다.

\textbf{검정력 분석.}
$n = 7$ 기능 영역이라는 작은 유효 표본 크기에서 방법 간 차이를 탐지하기 위한 통계적 검정력을 사후(post hoc)로 평가하였다.
대응표본(paired) 설계에서 유의 수준 $\alpha = 0.05$, $n = 7$일 때, 효과 크기 Cohen's $d \geq 1.4$에서 검정력 0.80이 달성된다.
MutClust-Hybrid와 MutClust-Fixed 간의 관찰된 hotspot-score 차이(0.778 -- 0.646 = 0.132)를 부트스트랩 표준 오차로 나눈 효과 크기는 $d \approx 2.1$로, 이 쌍별 비교가 충분한 검정력을 가짐을 확인한다.
다만 MutClust-Fixed와 MutClust-Original 간의 작은 차이(0.646 -- 0.638 = 0.008, $d \approx 0.15$)는 $n = 7$에서 탐지력이 극히 낮아, 이 두 방법의 동등성 주장에는 주의가 필요하다.

\textbf{Hotspot-score 가중치 민감도.}
Hotspot-score의 동일 가중치($1/3$) 설정의 강건성을 검증하기 위해, 171개의 가중치 조합($w_r + w_p + w_s = 1$, 단위 0.05)에서 방법 순위를 분석하였다(표~\ref{tab:weight_sensitivity}).
MutClust-Hybrid는 전체 가중치 공간의 \textbf{71.9\%}(123/171)에서 1위를 유지하였으며, HDBSCAN-Auto는 28.1\%(48/171)에서 1위를 차지하였다.
순위 전환은 recall 가중치가 0.60 이상이고 precision 가중치가 0.20 이하인 극단적 recall 중시 영역에서만 발생한다.

\begin{table}[htbp]
\centering
\caption{Hotspot-score 가중치 민감도 분석: 7가지 응용 시나리오에서의 방법 순위}
\label{tab:weight_sensitivity}
\begin{tabular}{lccccc}
\toprule
\textbf{시나리오} & \textbf{$(w_r, w_p, w_s)$} & \textbf{1위} & \textbf{1위 점수} & \textbf{2위} & \textbf{2위 점수} \\
\midrule
동일 가중치 (기본값) & (0.33, 0.33, 0.34) & Hybrid & 0.777 & Fixed & 0.644 \\
Recall 중시 & (0.50, 0.25, 0.25) & Hybrid & 0.748 & HDBSCAN & 0.689 \\
Precision 중시 & (0.25, 0.50, 0.25) & Hybrid & 0.826 & Fixed & 0.732 \\
Stability 중시 & (0.25, 0.25, 0.50) & Hybrid & 0.760 & HDBSCAN & 0.650 \\
\midrule
\textbf{조기경보 (감시)} & (0.60, 0.20, 0.20) & \textbf{HDBSCAN} & 0.740 & Hybrid & 0.731 \\
백신 설계 & (0.20, 0.60, 0.20) & Hybrid & 0.854 & Fixed & 0.784 \\
모니터링 & (0.20, 0.20, 0.60) & Hybrid & 0.749 & HDBSCAN & 0.678 \\
\bottomrule
\end{tabular}
\begin{minipage}{\textwidth}
\smallskip
\footnotesize 기본값 행의 가중치 $(0.33, 0.33, 0.34)$는 정확히 $1/3$이 아니므로, 표~\ref{tab:hotspot_scores}의 동일 가중치 결과(Hybrid 0.778, Fixed 0.646)와 소수점 셋째 자리에서 차이가 발생한다.
\end{minipage}
\end{table}

이 결과는 두 가지 중요한 함의를 가진다.
첫째, MutClust-Hybrid의 1위 순위는 가중치 선택에 대해 강건하여, 동일 가중치 설정이 특정 가중치 편향에 의한 인공적 결과가 아님을 확인한다.
둘째, 조기경보 목적(recall $\geq 0.60$)에서는 HDBSCAN-Auto가 더 적합할 수 있으며, 이는 MutBench 프레임워크가 단일 ``최적'' 방법을 지정하는 것이 아니라 \textbf{응용 맥락에 따른 방법 선택을 지원}하는 도구임을 보여준다.

쌍별 비교는 MutClust-Hybrid가 모든 다른 방법에 대해 유의하게 높은 hotspot-score를 달성함을 확인하였다(모든 비교에서 $p < 0.001$, 점수 차이 범위: MutClust-Fixed 대비 0.130에서 슬라이딩 윈도우 기준선 대비 0.811). 7개 방법의 쌍별 비교는 $\binom{7}{2} = 21$회의 다중 검정을 수반하며, Bonferroni 보정 후 유의 수준은 $\alpha_{\text{adj}} = 0.05 / 21 = 0.0024$이다. 모든 쌍별 $p < 0.001 < \alpha_{\text{adj}}$이므로 다중 비교 보정 후에도 결론이 유지된다.
이 결과는 표~\ref{tab:hotspot_scores}의 성능 순위가 ground truth 표본 추출 불확실성에 대해 견고하며 우연에 기인하지 않음을 확립한다.

\subsection{합성 데이터 설계의 순환성 검증}
\label{subsec:circularity_validation}

합성 벤치마크의 H-score 생성 과정에서 CCM 시드 위치(AA 484, 501, 417, 452)에 0.8--1.2 범위의 극단적 피크를 명시적으로 삽입하였으므로, CCM 기반 방법이 이 위치를 시드로 선택하여 부당한 이점을 얻을 가능성이 있다.
이 순환성 우려를 검증하기 위해, CCM 시드 피크를 제거하고 영역 수준 패턴(RBD 0.05--0.3, NTD 0.05--0.2, furin 0.1--0.5 등)만 유지한 ``블라인드'' 합성 데이터를 생성하여 동일한 벤치마크를 수행하였다.

블라인드 데이터에서 해당 위치의 H-score는 RBD 영역 배경 수준(0.09--0.25)으로 하락하였으나, 7개 방법의 F1 순위는 원본과 높은 상관을 보였다(Spearman $\rho = 0.821$, $p = 0.023$).
특히 MutClust-Hybrid는 원본(F1 = 0.785)과 블라인드(F1 = 0.786) 모두에서 1위를 유지하였고, F1 변화량은 $-$0.001로 사실상 동일하였다.
CCM 기반 3개 방법(Original, Fixed, dNdS)의 F1 변화도 0.002--0.004에 불과하여, 개별 시드 피크 제거가 탐지 결과에 미치는 영향이 무시할 수 있는 수준임을 확인하였다.
HDBSCAN-Auto와 OPTICS는 비CCM 방법으로서 순위 변동이 관찰되었으나(5$\leftrightarrow$6위), 이는 유사한 F1 값(0.137--0.150) 간의 미세 변동이다.

이 결과는 CCM 시드의 역할이 클러스터 ``시작점'' 제공에 한정되며, 최종 성능은 주변 영역의 밀도 구조에 의해 결정됨을 시사한다.
따라서 합성 벤치마크의 주요 결론---MutClust-Hybrid의 우위와 방법 간 성능 계층 구조---은 CCM 시드 순환성에 의한 인공물이 아니다.

%% ============================================================
\section{민감도 분석}
\label{sec:sensitivity_analysis}

\subsection{샘플 크기 민감도}

안정적인 핫스팟 탐지에 필요한 최소 입력 서열 수를 결정하기 위해, MutClust-Hybrid를 100에서 10,000 서열 범위의 샘플 크기에서 평가하였다(크기당 5회 반복).

H-score 순위 상관(Spearman $\rho$)은 부분추출과 전체 데이터셋 간에 $n = 1{,}000$에서 0.997을 초과하였으며 $n = 10{,}000$에서 1.000에 도달하였다.
Hotspot-score는 $n = 1{,}000$까지 약 0.62에서 안정화되었으며, 반복 간 분산이 최소(표준 편차 $< 0.015$)였다.
가장 작은 샘플 크기($n = 100$)에서도 hotspot-score는 0.52 이상을 유지하였고 Spearman $\rho > 0.94$를 보여, 제한된 데이터셋에서도 변이 강도의 순위 순서가 보존됨을 나타낸다.

이 결과는 약 1,000개의 서열이 안정적인 탐지 결과에 충분하며, 본 벤치마크 데이터셋의 10,000개 서열은 이 역치를 상당히 초과하는 여유를 제공함을 시사한다.

\subsection{파라미터 민감도}

샘플 크기에 대한 수렴 특성은 그림~\ref{fig:mutbench_sensitivity}에 시각화하였다.

\begin{figure}[htbp]
\centering
\includegraphics[width=0.9\textwidth]{sensitivity_heatmap.png}
\caption{MutClust 파라미터 민감도 히트맵. 64개 파라미터 조합($\gamma \times d \times$ MinPts)에 대한 전인자 격자 탐색 결과, F1 성능이 $\gamma$(엡실론 스케일링 인자)에 가장 민감하고, $d$(감쇠 인자)가 그 다음이며, MinPts의 영향은 미미함을 보여준다.}
\label{fig:mutbench_sensitivity}
\end{figure}

HDBSCAN 파라미터 민감도 분석에서는 16개 조합 전체에서 F1이 0.140--0.153 범위(변동폭 0.013)로 안정적이었으며, 탐지 위치 수는 약 12,400--14,900 범위에서 유지되었다(표~\ref{tab:hdbscan_sensitivity}).
이 파라미터 비민감성은 밀도 기반 접근법의 핵심 장점이며, MutClust-Hybrid가 경계 결정에 HDBSCAN을 사용하는 근거를 제공한다.

\begin{table}[htbp]
\centering
\caption{HDBSCAN 파라미터 민감도 분석}
\label{tab:hdbscan_sensitivity}
\begin{tabular}{rrrrrr}
\toprule
\textbf{min\_cluster\_size} & \textbf{min\_samples} & \textbf{정밀도} & \textbf{재현율} & \textbf{F1} & \textbf{탐지 수} \\
\midrule
3 & 2 & 0.083 & 0.869 & 0.152 & 13,054 \\
5 & 3 & 0.079 & 0.944 & 0.146 & 14,916 \\
10 & 5 & 0.081 & 0.962 & 0.149 & 14,939 \\
15 & 7 & 0.076 & 0.899 & 0.140 & 14,860 \\
\bottomrule
\end{tabular}
\end{table}

\subsection{dN/dS 역치 민감도}

dN/dS 필터의 역치 선택이 핫스팟 탐지 성능에 미치는 영향을 체계적으로 평가하기 위해, 6개의 역치(0.5, 0.8, 1.0, 1.5, 2.0, 3.0)에서 MutClust-dNdS와 MutClust-Hybrid를 합성 벤치마크에 적용하였다.
각 역치에서 dN/dS 필터 없이 클러스터를 탐지한 후, 사후적으로 dN/dS 주석을 적용하고 역치 이하의 클러스터를 제거하는 방식으로 분석하였다.
이 분석은 합성 데이터에서 참조 서열 대비 코돈 수준의 동의·비동의 치환을 인위적으로 계산할 수 있는 경우에만 적용되며, 합성 변이 생성 시 코돈 구조가 보존된 333개 변이 위치를 대상으로 하였다.

\begin{table}[htbp]
\centering
\caption{dN/dS 역치에 따른 핫스팟 탐지 성능 변화 (합성 데이터, 333개 합성 변이)}
\label{tab:dnds_threshold}
\small
\begin{tabular}{llrrrrrr}
\toprule
\textbf{방법} & \textbf{역치} & \textbf{정밀도} & \textbf{재현율} & \textbf{F1} & \textbf{유지} & \textbf{제외} & \textbf{HS} \\
\midrule
\multirow{7}{*}{MutClust-dNdS}
& 없음 & 0.988 & 0.506 & 0.669 & 30 & 0 & --- \\
& 0.5 & 0.985 & 0.414 & 0.583 & 22 & 8 & 0.580 \\
& 0.8 & 0.981 & 0.334 & 0.499 & 20 & 10 & 0.533 \\
& \textbf{1.0} & 0.976 & 0.261 & 0.412 & 17 & 13 & 0.502 \\
& 1.5 & 0.961 & 0.159 & 0.273 & 14 & 16 & 0.461 \\
& 2.0 & 0.952 & 0.112 & 0.200 & 11 & 19 & 0.429 \\
& 3.0 & 0.976 & 0.065 & 0.121 & 8 & 22 & 0.410 \\
\midrule
\multirow{7}{*}{MutClust-Hybrid}
& 없음 & 0.968 & 0.659 & 0.785 & 8 & 0 & --- \\
& 0.5 & 0.982 & 0.658 & 0.788 & 6 & 2 & 0.775 \\
& 0.8 & 0.982 & 0.658 & 0.788 & 6 & 2 & 0.770 \\
& \textbf{1.0} & 0.986 & 0.636 & 0.774 & 4 & 4 & 0.729 \\
& 1.5 & 0.973 & 0.057 & 0.107 & 2 & 6 & 0.453 \\
& 2.0 & 1.000 & 0.012 & 0.024 & 1 & 7 & 0.418 \\
& 3.0 & 0.000 & 0.000 & 0.000 & 0 & 8 & 0.000 \\
\bottomrule
\end{tabular}
\begin{minipage}{\textwidth}
\smallskip
\footnotesize 유지/제외: dN/dS 역치 적용 후 유지/제거된 클러스터 수. HS: hotspot-score. 기본 역치($> 1.0$)를 굵게 표시하였다.
\end{minipage}
\end{table}

표~\ref{tab:dnds_threshold}에 나타난 바와 같이, 두 방법 모두 역치 증가에 따라 재현율이 단조 감소하는 반면 정밀도는 비교적 안정적으로 유지된다.
MutClust-Hybrid는 전체 클러스터 수가 8개로 적기 때문에 역치 변화에 더 민감하며, 역치 1.0과 1.5 사이에서 성능이 급락한다(F1: 0.774 $\to$ 0.107).
역치 0.5--1.0 범위에서 MutClust-Hybrid의 성능이 안정적이며(F1: 0.774--0.788), 기본 역치 $> 1.0$(중립 진화 기대치)은 생물학적으로 해석 가능한 기준으로서 적절하다.
그러나 역치 $> 1.5$ 이상에서는 두 방법 모두 급격한 성능 저하를 보이므로, 엄격한 양성 선택 기준을 적용할 때는 주의가 필요하다.

%% ============================================================
\section{시간적 일반화 (GISAID 2024--2025)}
\label{sec:temporal}

MutBench 결과가 시간적으로 상이한 데이터에 대해 일반화되는지 평가하기 위해, 2024년 1월부터 2025년 12월까지 수집된 9,957개의 GISAID 전장 SARS-CoV-2 게놈 서열(숙주: 인간)을 분석하였으며, 시간적 대표성을 위해 주별로 부분추출하였다.

시간적 검증은 핵심적인 현상을 드러내었다: \textbf{H-score 포화}.
원래의 2020--2021 데이터셋은 114개 위치에서 비영 H-score를 산출하였으나(표준 편차 = 1.88), 2024--2025 데이터셋은 모든 1,273개 Spike 위치에서 비영 H-score를 산출하였다(평균 = 8.47, 표준 편차 = 0.098).
이 거의 균일한 분포는 위치 간 대비가 사실상 존재하지 않으므로 H-score의 판별력을 제거한다.

그 결과, 모든 방법이 약 0.33의 정밀도를 보이며, 이는 ground truth 위치의 기저율(1,251 / 3,819 = 0.327)과 거의 일치한다.
MutClust-Fixed와 MutClust-dNdS는 모든 3,819개 위치를 단일 클러스터로 탐지하였으며(정밀도 = 0.328, 재현율 = 1.0), MutClust-Hybrid는 일부 선택성을 유지하였으나(3,225개 위치, 377개 클러스터) 정밀도의 의미 있는 향상 없이 유사한 결과를 보였다(표~\ref{tab:temporal}).

\begin{table}[htbp]
\centering
\caption{GISAID 2024--2025 데이터를 이용한 시간적 검증 결과}
\label{tab:temporal}
\begin{tabular}{lrrrr}
\toprule
\textbf{방법} & \textbf{정밀도} & \textbf{재현율} & \textbf{탐지 수} & \textbf{클러스터} \\
\midrule
MutClust-Fixed & 0.328 & 1.000 & 3,819 & 1 \\
MutClust-dNdS & 0.328 & 1.000 & 3,819 & 1 \\
MutClust-Hybrid & $\sim$0.33 & $\sim$1.0 & 3,225 & 377 \\
\bottomrule
\end{tabular}
\begin{minipage}{\textwidth}
\smallskip
\footnotesize 모든 1,273개 Spike 위치에서 비영 H-score가 관찰되어 판별력이 제거되었다. 정밀도가 ground truth 기저율($\sim$0.33)에 수렴한다.
\end{minipage}
\end{table}

이 H-score 포화는 방법론적 인공물이 아닌 생물학적 현실이다: Omicron 시대 계통(BA.2.86, JN.1, KP.2, XEC)은 Wuhan-Hu-1 참조에 대해 광범위한 돌연변이를 축적하여, 위치 수준의 변이 빈도를 비정보적으로 만든다.
이 발견은 고정된 조상 참조에 대해 계산된 H-score가 바이러스 집단이 해당 참조로부터 분기함에 따라 판별력을 잃는다는 핵심 한계를 검증한다.

%% ============================================================
\section{다중 Ground Truth 평가}
\label{sec:multi_gt}

단일 ground truth로는 변이의 중요성의 모든 차원을 포착할 수 없다.
이를 해결하기 위해, 각각 서로 다른 생물학적 현상을 측정하는 세 가지 상보적 ground truth에 대해 탐지 방법을 평가하였다:

\begin{itemize}
\item \textbf{기능 영역}(439개 아미노산 위치, Spike 1,273 AA의 34.5\%): 알려진 도메인 주석(RBD, NTD 초과부위, 퓨린 절단, 융합 펩타이드, HR1, HR2)\footnote{본 분석에서 사용된 439 AA는 도메인 주석의 아미노산 수준 합집합으로, 합성 벤치마크에서 사용된 뉴클레오티드 수준 ground truth(1,251 nt, 중첩 제거 후 417 AA)보다 넓은 영역 경계 정의를 포함한다.}.
\item \textbf{DMS-escape}(252개 위치, 19.8\%): 실험적으로 측정된 단일 돌연변이 세포 진입 fitness 효과~\cite{dadonaite2023dms}(상위 20백분위).
\item \textbf{수렴 진화}(97개 위치, 7.6\%): 독립적 SARS-CoV-2 계통에서 반복적 양성 선택을 받은 문헌 기반 부위~\cite{cao2023ba286,carabelli2023convergent}. RBD 면역 회피 위치, ACE2 결합 강화 돌연변이, 퓨린 절단 부위 적응, S2 수렴 치환을 포함한다(이 확장 집합은 문헌에서 보고된 모든 수렴 진화 위치를 포함하며, 9-병원체 벤치마크에서 사용된 엄격한 25개 위치 Layer~A 집합과 구분된다).
\end{itemize}

세 ground truth는 대체로 비중첩적이다: DMS-escape와 수렴 진화 위치 간의 Jaccard 유사도는 0.006으로, 거의 완전한 독립성을 나타낸다.
이는 집단 빈도, 실험적 fitness, 진화적 수렴이 변이 중요성의 근본적으로 다른 측면을 포착함을 확인한다.

결정적으로, 최우수 방법은 ground truth에 따라 달랐다(표~\ref{tab:multi_gt}).

\begin{table}[htbp]
\centering
\caption{다중 ground truth에 대한 방법 비교}
\label{tab:multi_gt}
\begin{tabular}{llrrr}
\toprule
\textbf{Ground Truth} & \textbf{최우수 방법} & \textbf{F1} & \textbf{농축비} & \textbf{위치 수} \\
\midrule
기능 영역 (439) & TC-freq-v2 & 0.474 & --- & 439 \\
DMS-escape (252) & Saturated H-score & 0.315 & --- & 252 \\
수렴 진화 (97) & CH-score & 0.188 & 1.70$\times$ & 97 \\
수렴 진화-엄격$^*$ & TC-freq-v3 & 0.108 & 3.43$\times$ & --- \\
\bottomrule
\end{tabular}
\begin{minipage}{\textwidth}
\smallskip
\footnotesize $^*$다중 출처에서 확인된 엄격한 수렴 진화 부위만 포함. Jaccard(DMS, 수렴 진화) = 0.006으로 세 ground truth의 독립성이 확인된다.
\end{minipage}
\end{table}

기능 영역에서는 TC-freq-v2가 최고 F1(0.474)을 달성하였고, DMS-escape에서는 포화 H-score가 최고 성능(F1 = 0.315)을 보였으며, 수렴 진화에서는 CH-score가 최고 F1(0.188, 농축비 1.70배)을 달성하였다.
엄격한 수렴 진화 부분집합(다중 출처 위치만)에서는 TC-freq-v3가 최고 농축비(3.43배, F1 = 0.108)를 달성하였다.

이러한 상이한 순위는 어떤 단일 탐지 방법도 핫스팟 중요성의 모든 생물학적 차원에서 지배적이지 않음을 입증한다.
이 다중 ground truth 분석은 핵심적인 방법론적 통찰을 제공한다: 집단 빈도와 실험적 fitness는 변이 중요성의 상보적이지 중복적이지 않은 측정이다.
단일 ground truth에 의존하는 벤치마킹 프레임워크는 해당 특정 생물학적 차원에 최적화된 방법에 대한 체계적 편향의 위험을 감수하면서, 기능적 중요성의 직교적 측면을 포착하는 방법을 불이익으로 처리할 수 있다.

\section{실제 GISAID 2020--2021 H-score 벤치마크}
\label{sec:real_gisaid_benchmark}

합성 데이터의 순환성 우려를 해소하기 위해, MutClust 논문~\cite{youn2025mutclust}에서 사용된 실제 GISAID 2020--2021 H-score 데이터를 입력으로 사용한 벤치마크를 수행하였다.
이 데이터는 2020년 1월부터 2021년 12월까지 수집된 SARS-CoV-2 서열에서 계산된 1,273개 Spike 아미노산 위치별 H-score이며, 합성 생성이 아닌 실제 관측 데이터이다.
뉴클레오티드 수준으로 매핑한 결과, 29,903개 게놈 위치 중 342개(1.14\%)만이 비영 H-score를 가져, 합성 데이터의 균일한 배경 잡음과 질적으로 다른 희소(sparse) 분포를 나타낸다.

결과는 표~\ref{tab:real_gisaid}에 제시하였다.

\begin{table}[htbp]
\centering
\caption{실제 GISAID 2020--2021 H-score 데이터에 대한 방법 비교 (342개 비영 위치, 1,251개 ground truth 뉴클레오티드 위치)}
\label{tab:real_gisaid}
\begin{tabular}{lrrrrrr}
\toprule
\textbf{방법} & \textbf{정밀도} & \textbf{재현율} & \textbf{F1} & \textbf{농축비} & \textbf{안정성} & \textbf{Hotspot-score} \\
\midrule
MutClust-Original & 0.333 & 0.012 & 0.023 & 7.97 & 0.372 & 0.239 \\
MutClust-Fixed & 0.295 & 0.079 & 0.125 & 7.04 & 0.784 & 0.386 \\
MutClust-dNdS & 0.295 & 0.079 & 0.125 & 7.04 & 0.784 & 0.386 \\
MutClust-Hybrid & 0.305 & 0.077 & 0.123 & 7.28 & 0.766 & 0.383 \\
HDBSCAN-Auto & --- & --- & --- & --- & --- & --- \\
\midrule
빈도 기준선 & 0.289 & 0.079 & 0.124 & 6.92 & --- & --- \\
\textbf{슬라이딩 윈도우} & \textbf{0.323} & \textbf{0.600} & \textbf{0.420} & \textbf{7.73} & --- & --- \\
\bottomrule
\end{tabular}
\begin{minipage}{\textwidth}
\smallskip
\footnotesize HDBSCAN-Auto는 비영 위치 수(342개)가 불충분하여 클러스터를 형성하지 못하였다. 슬라이딩 윈도우가 넓은 탐지 창(20위치)을 통해 최고 F1(0.420)을 달성하였다.
\end{minipage}
\end{table}

합성 벤치마크와 비교하여 세 가지 핵심 차이가 관찰된다.

\textbf{첫째, 절대 성능의 현저한 하락.}
모든 MutClust 변이체의 F1이 합성 대비 0.54--0.66 하락하였다(예: MutClust-Hybrid: 0.785 $\to$ 0.123).
이는 실제 데이터의 희소성에 기인한다: ground truth 1,251개 위치 중 상당수(HR1, HR2 등 보존 영역)는 H-score = 0이며, H-score가 0인 위치는 어떤 밀도 기반 방법으로도 탐지할 수 없다.
이 결과는 합성 벤치마크가 실제 성능을 과대추정함을 정직하게 드러내며, 밀도 기반 핫스팟 탐지가 ``변이하는 기능 부위''만 탐지 가능하고 ``보존된 기능 부위''는 원리적으로 탐지 불가능하다는 방법론적 한계를 확인한다.

\textbf{둘째, 방법 순위의 변화.}
합성 데이터에서 MutClust-Hybrid가 최고 hotspot-score(0.778)였으나, 실제 데이터에서는 MutClust-Fixed(0.386)와 거의 동일한 성능을 보였다(0.383).
HDBSCAN 경계 결정의 이점은 연속적 H-score 분포(합성 데이터)에서 발현되지만, 이산적이고 희소한 분포(실제 데이터)에서는 CCM 시드 탐지가 경계 결정보다 지배적이다.
슬라이딩 윈도우가 최고 F1(0.420)을 달성한 것은 넓은 탐지 창이 희소 데이터에서 인접 기능 영역을 포괄하기 유리함을 나타낸다.

\textbf{셋째, 농축비의 일관적 유의성.}
절대 F1은 낮지만, 모든 방법의 농축비가 6.9--8.0배로 일관되게 높아(무작위 기대 대비), 탐지된 위치가 기능 영역에 유의하게 농축됨은 합성·실제 데이터 모두에서 확인된다.
이는 H-score 기반 탐지의 생물학적 타당성이 데이터 유형에 관계없이 유지됨을 의미한다.

이 결과는 MutBench의 핵심 방법론적 교훈을 강화한다: 방법의 상대적 우위는 입력 데이터의 H-score 분포 특성---연속적 vs. 희소적, 균일 배경 vs. 이산적 배경---에 따라 달라지며, 합성 벤치마크의 결론을 실제 데이터에 무조건 외삽해서는 안 된다.

\section{외부 독립 방법과의 비교}
\label{sec:external_comparison}

MutClust 변이체 간의 내부 비교를 보완하기 위해, 두 가지 독립적인 클러스터링 기반 방법---OPTICS와 커널 밀도 추정(KDE)---을 동일한 합성 벤치마크에서 평가하였다.
OPTICS는 DBSCAN과 유사한 밀도 기반 클러스터링이나, 도달 가능성 거리(reachability distance)를 통해 가변 밀도 클러스터를 자동으로 추출한다(min\_samples=5, $\xi$=0.05).
KDE는 H-score의 커널 밀도 함수에서 90번째 백분위수 이상의 피크를 탐지한 후 인접 위치를 클러스터로 확장하는 비모수적 접근법이다.

\begin{table}[htbp]
\centering
\caption{외부 독립 방법을 포함한 벤치마크 결과. 합성 H-score 데이터, 7개 Spike 기능 영역 ground truth 사용.}
\label{tab:external_methods}
\small
\begin{tabular}{lcccccc}
\toprule
방법 & 정밀도 & 재현율 & F1 & 안정성 & Hotspot-score \\
\midrule
MutClust-Hybrid & 0.968 & 0.659 & 0.785 & 0.706 & \textbf{0.778} \\
MutClust-Original & 0.991 & 0.504 & 0.668 & 0.418 & 0.638 \\
MutClust-Fixed & 0.988 & 0.506 & 0.669 & 0.445 & 0.646 \\
MutClust-dNdS & 0.988 & 0.506 & 0.669 & 0.445 & 0.646 \\
HDBSCAN-Auto & 0.079 & 0.944 & 0.146 & 0.788 & 0.604 \\
KDE & \textbf{1.000} & 0.049 & 0.093 & 0.622 & 0.557 \\
OPTICS & 0.087 & 0.546 & 0.150 & 0.473 & 0.369 \\
\bottomrule
\end{tabular}
\end{table}

표~\ref{tab:external_methods}에 나타난 바와 같이, MutClust-Hybrid는 외부 독립 방법을 포함한 비교에서도 가장 높은 hotspot-score(0.778)를 달성하였다.
KDE는 완벽한 정밀도(1.000)를 보였으나 극단적으로 보수적이어서 61개 위치만을 탐지하였으며(재현율 0.049), OPTICS는 997개 클러스터에 7,847개 위치를 탐지하였으나 정밀도가 0.087에 불과하였다.
이 결과는 도메인 특화 사전 지식(CCM 시드)과 데이터 기반 경계 결정(HDBSCAN)의 결합이 범용 클러스터링 방법보다 우월함을 독립적으로 확인한다.

\textbf{OPTICS/KDE 파라미터 민감도.}
외부 방법의 파라미터 설정이 결과에 미치는 영향을 배제하기 위해, 합성 데이터에서 OPTICS(min\_samples $\in \{2, 3, 5, 10, 20, 50\}$, $\xi \in \{0.01, 0.03, 0.05, 0.1, 0.2\}$)와 KDE(bandwidth $\in \{1, 3, 5, 10, 20, 50\}$, 역치 백분위 $\in \{70, 80, 85, 90, 95, 99\}$)의 파라미터 민감도 분석을 수행하였다.
OPTICS는 30개 조합에서 F1이 0.022--0.707 범위로 넓은 파라미터 의존성을 보였다.
최적 파라미터(min\_samples=20, $\xi$=0.10)에서 OPTICS는 F1=0.707(정밀도 0.563, 재현율 0.949)을 달성하였으나, 이는 MutClust-Hybrid(F1=0.785, 정밀도 0.968)에 비해 F1이 0.078 낮고 정밀도가 0.405 낮아, CCM 시드 기반 특이성의 이점이 확인된다.
KDE는 36개 조합 전체에서 F1=0.000으로, bandwidth와 역치 백분위에 관계없이 유효한 탐지를 수행하지 못하였다.
실제 GISAID 데이터에서는 비영 위치의 희소성(342개)으로 인해 OPTICS가 3개 위치만을 탐지하였고 KDE는 유효한 클러스터를 형성하지 못하여, 희소 분포에서의 범용 클러스터링 방법의 추가적 한계가 확인되었다.
이 결과는 MutClust-Hybrid의 우위가 특정 파라미터 설정에 대한 OPTICS/KDE의 불리한 비교가 아닌, 도메인 특이적 시드 탐지의 구조적 이점에 기인함을 확인한다.

\section{교차 병원체 검증: Influenza A H3N2 HA}
\label{sec:influenza_validation}

MutBench의 교차 병원체 일반화 가능성을 평가하기 위해, Influenza A H3N2의 헤마글루티닌(HA) 유전자에 동일한 벤치마크 파이프라인을 적용하였다.
NCBI Influenza Virus Resource에서 1,967개의 인간 유래 H3N2 HA 단백질 서열(540--580 AA)을 수집하고, 합의 서열 대비 위치별 변이 빈도와 Shannon 엔트로피로부터 H-score를 계산하였다.

Ground truth로는 H3N2 HA의 5개 항원 부위(antigenic sites A--E)를 사용하였으며~\cite{wiley1981structural,wilson1990structural}, 이는 수십 년의 혈청학적·구조적 연구에 의해 정의된 확립된 기능 영역이다(101개 AA 위치, 303개 뉴클레오티드 위치). Smith et al.~\cite{smith2004antigenic}은 항원 지도작성을 통해 이 부위들의 소수 치환이 주요 항원 전환(antigenic drift)을 주도함을 입증하였다.

\begin{table}[htbp]
\centering
\caption{Influenza A H3N2 HA 교차 병원체 벤치마크 결과 (1,967개 서열, 5개 항원 부위 ground truth)}
\label{tab:influenza}
\begin{tabular}{lrrrrrr}
\toprule
\textbf{방법} & \textbf{정밀도} & \textbf{재현율} & \textbf{F1} & \textbf{농축비} & \textbf{안정성} & \textbf{Hotspot-score} \\
\midrule
MutClust-Original & 0.182 & 1.000 & 0.309 & 1.02 & 0.535 & 0.572 \\
\textbf{MutClust-Fixed} & \textbf{0.183} & \textbf{1.000} & \textbf{0.309} & \textbf{1.02} & \textbf{0.799} & \textbf{0.660} \\
MutClust-Hybrid & 0.178 & 0.812 & 0.292 & 1.00 & 0.742 & 0.577 \\
HDBSCAN-Auto & 0.155 & 0.366 & 0.218 & 0.87 & 0.750 & 0.424 \\
KDE & 0.410 & 0.083 & 0.138 & 2.30 & 0.605 & 0.366 \\
\midrule
빈도 기준선 & 0.421 & 0.238 & 0.304 & 2.36 & --- & --- \\
\textbf{슬라이딩 윈도우} & \textbf{0.402} & \textbf{0.462} & \textbf{0.430} & \textbf{2.25} & --- & --- \\
\bottomrule
\end{tabular}
\end{table}

H3N2 결과는 SARS-CoV-2와 세 가지 핵심 차이를 보인다.

\textbf{첫째, 방법 순위의 전환.}
SARS-CoV-2 합성 벤치마크에서 최고 성능이었던 MutClust-Hybrid(hotspot-score 0.778)가 H3N2에서는 MutClust-Fixed(0.660)에 뒤처졌다(0.577).
이는 HA 유전자의 H-score 분포 특성에 기인한다: 566개 위치 모두가 비영 H-score를 가져(평균 1.233), CCM 시드가 거의 모든 위치를 후보로 선정하며, HDBSCAN 경계 결정이 추가적인 구별력을 제공하지 못한다.
이 결과는 ``최적 방법은 입력 데이터의 H-score 분포에 의존한다''는 실제 GISAID 벤치마크(제~\ref{sec:real_gisaid_benchmark}절)의 교훈을 교차 병원체 수준에서 재확인한다.

\textbf{둘째, 농축비의 현저한 감소.}
SARS-CoV-2에서 1.9--23.7배였던 농축비가 H3N2에서는 0.87--2.36배로 낮아졌다.
이는 HA의 짧은 유전자 길이(1,698 nt vs. 3,822 nt)와 높은 전반적 변이율(모든 위치 비영)이 결합되어, 항원 부위와 비항원 부위 간의 변이 빈도 대비가 약하기 때문이다.
HA의 항원 부위는 양성 선택 하에 있으나, 비항원 부위도 유전적 부동(genetic drift)에 의한 상당한 변이를 축적하여 배경 대비 신호가 약화된다.

\textbf{셋째, 슬라이딩 윈도우의 우수 성능.}
F1 기준으로 슬라이딩 윈도우(0.430)가 모든 밀도 기반 방법(0.138--0.309)을 능가하였다.
이는 HA 항원 부위가 비교적 작고 컴팩트한 영역(평균 20 AA)으로, 고정 창 크기의 슬라이딩 윈도우가 이러한 구조에 적합함을 시사한다.
반면 SARS-CoV-2 Spike의 기능 영역은 RBD(223 AA)처럼 크고 이질적인 구조를 포함하여 적응형 클러스터링이 유리하다.

이 교차 병원체 검증은 MutBench 프레임워크의 두 가지 핵심 가치를 입증한다:
(1)~동일한 벤치마크 파이프라인이 코드 변경 없이 다른 병원체에 적용 가능하며,
(2)~방법 순위의 병원체 의존성이 체계적 비교 없이는 발견될 수 없는 중요한 통찰이다.

\subsubsection{교차 병원체 비교}

표~\ref{tab:cross_pathogen}은 SARS-CoV-2와 H3N2에 대한 방법 순위를 종합적으로 비교한다.

\begin{table}[htbp]
\centering
\caption{2-병원체 교차 비교: 최우수 방법의 병원체 의존성}
\label{tab:cross_pathogen}
\small
\begin{tabular}{llrr}
\toprule
& & \multicolumn{2}{c}{\textbf{Hotspot-score}} \\
\cmidrule(lr){3-4}
\textbf{방법} & \textbf{유형} & \textbf{SARS-CoV-2} & \textbf{H3N2} \\
\midrule
MutClust-Original & CCM & 0.638 & 0.572 \\
MutClust-Fixed & CCM & 0.646 & \textbf{0.660} \\
MutClust-Hybrid & CCM+HDBSCAN & \textbf{0.778} & 0.577 \\
HDBSCAN-Auto & 밀도 기반 & 0.604 & 0.424 \\
\bottomrule
\end{tabular}
\begin{minipage}{\textwidth}
\smallskip
\footnotesize 각 병원체별 최고 hotspot-score를 \textbf{굵게} 표시. 최우수 방법이 두 병원체에서 상이하다. H3N2 벤치마크는 NCBI에서 수집한 실제 서열 기반 H-score를 사용하였다.
\end{minipage}
\end{table}

두 병원체에서 최우수 방법이 상이하다---SARS-CoV-2에서는 MutClust-Hybrid(0.778), H3N2에서는 MutClust-Fixed(0.660).
이 결과는 ``보편적 최우수 방법은 존재하지 않으며, 방법 선택은 입력 데이터의 특성(유전자 길이, H-score 분포, 신호-배경 대비)에 의존한다''는 MutBench의 핵심 교훈을 독립 병원체에서 확인한다.

생물학적 관점에서 이 패턴은 명확하게 해석된다: (1)~긴 유전자(Spike, 3,822 nt)에서 이질적 기능 영역이 분포할 때 적응형 하이브리드 접근(MutClust-Hybrid)이 유리하고, (2)~모든 위치가 높은 변이율을 보이는 짧은 유전자(HA, 1,698 nt)에서는 안정적 CCM 탐지(MutClust-Fixed)가 우수하다.
이는 H-score 분포의 형태(연속적 vs. 희소)와 유전자 길이가 최적 탐지 전략을 결정하는 핵심 요인임을 시사한다.

\section{3D 구조 분석}
\label{sec:structure_analysis}

선형 서열 기반 핫스팟 탐지의 구조적 타당성을 검증하기 위해, PDB 6VSB (Spike 단백질 prefusion 구조, Wrapp et al. 2020)의 C$\alpha$ 좌표를 이용하여 탐지된 핫스팟 위치의 3D 공간적 분포를 분석하였다.
핫스팟 위치 쌍의 8\,\AA\ 이내 접촉 비율은 5.80\%로, 무작위 위치 쌍(0.98\%)에 비해 \textbf{5.92배 높은 공간적 접촉 농축}을 보였다($z$ = 13.34, 1,000회 무작위 순열 검정 기준).
이는 선형 서열에서 탐지된 핫스팟이 3D 구조에서도 공간적으로 클러스터링됨을 정량적으로 입증하며, 탐지 결과의 생물학적 타당성을 뒷받침한다.

\section{계통발생 비독립성 시뮬레이션}
\label{sec:phylo_simulation}

H-score가 계통발생적 구조에 의해 편향될 수 있는지를 정량적으로 평가하기 위해, 합체 과정 기반 시뮬레이션을 수행하였다.

\subsection{계통발생 보정의 개념 증명}
\label{subsec:phylo_correction_poc}

합체 시뮬레이션(100회 반복, $N = 500$)에서 단일 기원의 ``창시자 변이''(빈도 90\%, D614G 유사)와 15회 독립 발생하는 ``수렴 변이''(E484K 유사), 50개의 중립 변이를 배치하였다.
H-score는 100\%의 반복에서 창시자 변이를 수렴 변이보다 높게 순위한 반면, Fitch 절약법 기반 homoplasy 횟수는 100\%에서 수렴 변이를 올바르게 최상위로 배치하였다.
H-score와 독립 발생 횟수 간 Spearman 순위 상관은 $\rho = -0.876$ ($p = 1.28 \times 10^{-64}$)으로, \textbf{빈도 기반 점수가 창시자 효과를 양성 선택과 체계적으로 혼동함}을 정량적으로 입증하였다.

MutBench ground truth에서 D614G(위치 614)는 단일 창시자 이벤트~\cite{korber2020tracking}에 의해 고정되었으므로 Layer~A(수렴 진화) 및 Layer~B(정제 선택) 모두에서 의도적으로 제외되었다.
이 결과는 계통발생 기반 독립 발생 횟수 추정이 빈도 기반 접근법의 창시자 효과 편향을 효과적으로 보정할 수 있음을 개념적으로 입증한다.
다만, 본 시뮬레이션은 이상적 조건($N = 500$, 완전 해결 계통수)이며, GISAID 규모($> 16$M 서열)에서는 polytomies, 계산 비용, 재조합 등의 추가적 도전이 존재한다.

\subsection{TreeTime ML 조상 재구성 검증}
\label{subsec:treetime_ml_validation}

TreeTime~\cite{sagulenko2018treetime}의 ML 기반 조상 재구성을 통합한 확장 검증($N = 2{,}000$, 50회 반복)에서, Fitch 절약법과 TreeTime-ML 모두 수렴 변이 대 비수렴 변이 판별에서 AUC = 1.000을 달성한 반면, H-score는 AUC = 0.000을 기록하였다.
$N = 500$에서 $N = 2{,}000$으로 스케일을 확장하여도 100\% 판별 정확도가 유지되었다.

\subsection{실제 GISAID 서열에서의 TreeTime 조상 재구성 검증}
\label{subsec:treetime_gisaid_validation}

합성 시뮬레이션을 넘어 실제 데이터에서 검증하기 위해, 500개 GISAID 서열(2024--2025)에 TreeTime ML 조상 재구성을 적용하였다.
전체 1,273개 Spike 위치에 대한 Spearman 순위 상관은 $\rho = -0.107$ ($p = 1.36 \times 10^{-4}$)로, H-score와 homoplasy 점수 간에 \textbf{약한 음의 상관}이 관찰되었다.
이는 두 점수 체계가 근본적으로 다른 진화적 신호를 포착함을 실증한다: H-score는 참조 대비 빈도를, homoplasy는 독립적 반복 발생을 반영한다.
핵심 위치에서 이 괴리가 명확하다: P681은 H-score 순위 2위이나 homoplasy 순위 1,134위(founder effect)이며, L452와 K417은 H-score = 0이나 homoplasy 91회(수렴 진화)를 보인다.
이 결과는 시뮬레이션에서 예측한 H-score의 founder effect 편향이 실제 GISAID 서열에서도 경험적으로 확인됨을 보여준다.

\subsection{Homoplasy 기반 스코어링 민감도 분석}
\label{subsec:homoplasy_sensitivity}

이 분석을 실제 계통발생 데이터로 검증하기 위해, 2020-H1 NCBI 데이터셋에서 서브샘플링한 500개 SARS-CoV-2 Spike 단백질 서열에 대해 TreeTime 최대우도 조상 재구성을 사용한 민감도 실험을 수행하였다.
FastTree(WAG 모델)로 계통수를 추론하고, TreeTime이 모든 내부 노드에서 조상 서열을 재구성하였으며, 1,273개 Spike 아미노산 위치 각각에 매핑된 독립적 변이 이벤트 수를 합산하여 위치별 homoplasy 횟수를 추출하였다.
이후 5개 알고리즘 패밀리의 11개 탐지 방법을 homoplasy 횟수 배열에 적용하고, Layer~A ground truth(25개 수렴 진화 위치)에 대해 MCC를 산출하였다.

\begin{table}[htbp]
\centering
\caption{SARS-CoV-2 Spike에서 H-score 대 TreeTime ML 추론 homoplasy 횟수의 탐지 방법별 MCC 비교 (500개 서열, 2020-H1).}
\label{tab:phylo_sensitivity}
\small
\begin{tabular}{lccccc}
\toprule
탐지 방법 & \multicolumn{2}{c}{H-score} & \multicolumn{2}{c}{TreeTime Homoplasy} & $\Delta$MCC \\
\cmidrule(lr){2-3} \cmidrule(lr){4-5}
 & 최적 공식 & MCC & MCC & 정밀도 & \\
\midrule
FreqThresh(p=85)     & rank(P*E) & 0.015 & $-0.036$ & 0.000 & $-0.051$ \\
FreqThresh(p=90)     & P*E       & 0.047 & $-0.029$ & 0.000 & $-0.076$ \\
FreqThresh(p=95)     & E*rare    & 0.076 & $-0.015$ & 0.000 & $-0.091$ \\
AdaptiveKnee         & E*rare    & 0.116 & $-0.033$ & 0.000 & $-0.149$ \\
LocalContrast(z=1.5) & minority\_E & 0.068 & $-0.033$ & 0.000 & $-0.101$ \\
Wavelet(t=1.5)       & E*rare    & 0.086 & $-0.016$ & 0.000 & $-0.102$ \\
ScoreDBSCAN(ep=10)   & rank(P*E) & 0.023 & $+0.033$ & 0.025 & $+0.010$ \\
\midrule
평균                 &           & 0.062 & $-0.018$ &       & $-0.080$ \\
\bottomrule
\end{tabular}
\end{table}

\emph{실제} TreeTime ML 추론 homoplasy 횟수를 사용한 표~\ref{tab:phylo_sensitivity}에서 세 가지 발견이 도출된다.
첫째, ScoreDBSCAN을 제외한 모든 탐지 방법이 음의 MCC(평균 $-0.018$)를 산출하며, 대부분 Layer~A ground truth에 대해 정밀도와 재현율이 0이다.
이 결과는 시간적 불일치로 설명된다: 2020-H1 서열은 우려 변이(Variants of Concern)의 출현보다 이전 시기이므로(Alpha, Delta, Omicron은 2020년 후반부터 출현), Layer~A ground truth의 수렴 진화 위치(예: N501Y, E484K, L452R)가 초기 팬데믹에서 아직 반복 변이로 발생하지 않았다.
최고 homoplasy는 257--320 위치(NTD 영역)에 집중되어, ground truth를 정의하는 RBD 부위의 적응적 수렴 진화보다는 초기 계통에서의 반복적 독립 결실을 반영하였다.

둘째, 창시자 효과 위치인 D614G는 homoplasy 횟수 0을 정확히 기록하여(독립적 출현이 아닌 단일 기원 확산에 부합), P681H는 1을 기록하여, TreeTime 재구성이 창시자 주도 고정과 수렴 진화를 정확히 구분함을 확인하였다.

셋째, H-score와 TreeTime homoplasy에서의 방법 순위 간 Spearman 순위 상관이 $\rho = 0.036$ ($p = 0.939$)으로, 사실상 상관이 없다.
이는 방법 순위가 스코어링 체계와 입력 데이터의 시간적 범위에 조건적임을 보여주며, 제~\ref{subsec:pahd_prototype}절에서 제안하는 병원체 적응적 접근법의 필요성을 강화한다.
9-병원체 벤치마크의 질적 결론---병원체마다 최적 방법이 다르다---은 스코어링 체계와 무관하게 유지될 것으로 기대되지만, 구체적인 양적 순위는 H-score 기반 스코어링과 입력 서열의 시간적 윈도우에 조건적이다.

\section{9-병원체 벤치마크 결과}
\label{sec:9pathogen_results}

제~\ref{sec:9pathogen_benchmark}절에서 제시한 9-병원체 확장 벤치마크를 적용하여, 9가지 스코어링 공식과 41개 탐지 방법의 조합을 9종 병원체에 걸쳐 체계적으로 평가하였다.
본 절에서는 벤치마크 규모, 병원체별 최적 방법 분석, 분산분석, 교차 검증, 통계 검정 결과를 순차적으로 보고한다.

\subsection{벤치마크 규모 및 데이터}
\label{subsec:9pathogen_scale}

총 3,321개의 평가(9개 스코어링 $\times$ 41개 탐지 방법 $\times$ 9개 병원체)를 수행하였다.
이 중 777개(23.4\%)의 평가에서 탐지 방법이 0개의 위치를 탐지하였으며(n\_detected = 0), 이는 SWAN 및 ScoreDBSCAN 일부 변이체가 특정 병원체의 점수 분포에서 $z$-정규화 역치 또는 밀도 기준을 충족하는 위치를 찾지 못하기 때문이다.
이러한 영탐지(zero-detection) 평가는 MCC = 0으로 처리된다. 단, ANOVA 분석에서는 영탐지 평가를 제외한 2,544개의 유효 평가를 사용하였으며, Type III 제곱합을 적용하였다.
각 평가는 3계층 ground truth(Layer A: 수렴 진화 위치, Layer B: 구속 영역, Layer C: DMS)에 대해 MCC, F1, 정밀도, 재현율, 농축비, 구속 영역 FPR을 산출한다.
DMS F1은 실험적 fitness 데이터가 이용 가능한 4종 병원체(SARS-CoV-2, H3N2, HIV-1, RSV)에서만 산출되었다.
모든 FASTA 서열은 NCBI GenBank에서 실제 accession을 통해 수집하였다.

\subsection{병원체별 최적 방법 분석}
\label{subsec:9pathogen_best_combos}

9종 병원체 각각에서 MCC를 기준으로 최적 스코어링--탐지 조합을 선정한 결과, \textbf{9/9 병원체에서 모두 상이한 최적 조합}이 관찰되었다(표~\ref{tab:9pathogen_best}).

\begin{table}[htbp]
\centering
\caption{9종 병원체별 최적 스코어링--탐지 조합 (MCC 기준)}
\label{tab:9pathogen_best}
\small
\begin{tabular}{llll}
\toprule
\textbf{병원체} & \textbf{최적 스코어링} & \textbf{최적 탐지 방법} & \textbf{MCC} \\
\midrule
SARS-CoV-2 & E*rare & CUSUM ($\delta$=0.5, $h$=4.0) & 0.118 \\
H3N2 & E\_only & SlidingTest ($p$=0.05) & 0.549 \\
Norovirus & P*E*rare & GradPeak ($p$=0.5) & 0.438 \\
HIV-1 & E*rare & SlidingTest ($p$=0.01) & 0.182 \\
Dengue & E*rare & FreqThresh ($p$=80) & 0.256 \\
RSV & P*E & FreqThresh ($p$=95) & 0.249 \\
Influenza B & H-score & Wavelet ($t$=2.0) & 0.361 \\
MERS & rank(P*E) & SlidingTest ($p$=0.05) & 0.088 \\
HCV & E\_only & CUSUM ($\delta$=0.5, $h$=5.0) & 0.443 \\
\bottomrule
\end{tabular}
\begin{minipage}{\textwidth}
\smallskip
\footnotesize 9개 병원체에서 9개의 고유한 (스코어링, 탐지) 조합이 최적으로 선정되었으며, 단일 조합이 두 개 이상의 병원체에서 동시에 최적인 경우는 없었다.
\end{minipage}
\end{table}

이 결과의 핵심 관찰은 다음과 같다.
첫째, 최적 스코어링 공식이 병원체마다 상이하다: E*rare가 SARS-CoV-2, HIV-1, Dengue에서 최적이나, H3N2와 HCV에서는 E\_only가, Influenza B에서는 H-score가, RSV에서는 P*E가, MERS에서는 rank(P*E)가 최적이다.
둘째, 최적 탐지 패밀리도 병원체마다 상이하다: SlidingTest(3회), CUSUM(2회), FreqThresh(2회), Wavelet(1회), GradPeak(1회)으로 분산된다.
셋째, MCC 범위가 0.088(MERS)에서 0.549(H3N2)까지 큰 변동을 보이며, 이는 병원체 간 변이 패턴의 구조적 차이(엔트로피 분포, 보존 영역 비율, 양성 선택 부위 수)를 반영한다.

\subsection{삼원분산분석 (Three-way ANOVA, $\omega^2$)}
\label{subsec:9pathogen_anova}

성능 변동의 주요 원천을 식별하기 위해, MCC를 종속변수로 하는 삼원분산분석을 수행하였다.
편향 보정된 효과 크기 $\omega^2$(식~\ref{eq:omega_squared})의 결과를 표~\ref{tab:9pathogen_anova}에 제시한다.

\begin{table}[htbp]
\centering
\caption{삼원분산분석 결과: MCC에 대한 각 요인의 $\omega^2$ 효과 크기}
\label{tab:9pathogen_anova}
\begin{tabular}{llr}
\toprule
\textbf{요인} & \textbf{유형} & \textbf{$\omega^2$} \\
\midrule
pathogen & 주효과 & 0.260 \\
family $\times$ pathogen & 이원 상호작용 & \textbf{0.285} \\
score $\times$ pathogen & 이원 상호작용 & 0.083 \\
score & 주효과 & 0.027 \\
family & 주효과 & 0.025 \\
score $\times$ family & 이원 상호작용 & 0.013 \\
\bottomrule
\end{tabular}
\begin{minipage}{\textwidth}
\smallskip
\footnotesize family: 탐지 방법 패밀리 (15수준), score: 스코어링 공식 (9수준), pathogen: 병원체 (9수준). $\omega^2$는 전체 분산 중 해당 요인이 설명하는 비율의 편향 보정 추정치이다.
\end{minipage}
\end{table}

\textbf{핵심 발견}: 탐지 패밀리 $\times$ 병원체 상호작용($\omega^2 = 0.285$)이 \textbf{모든 주효과 및 다른 상호작용보다 큰 최대 분산 원천}이다. Cohen의 효과 크기 기준~\cite{cohen1988statistical}에 따르면 $\omega^2 > 0.14$는 대(large) 효과에 해당하므로, 이 상호작용은 실질적으로 큰 효과이다.
이는 최적 탐지 방법이 병원체에 따라 체계적으로 달라지며, 병원체 독립적인 단일 최적 방법이 존재하지 않음을 의미한다.
병원체 주효과($\omega^2 = 0.260$)도 크며, 이는 병원체 간 변이 구조의 본질적 차이(엔트로피 분포 형태, 양성 선택 부위 밀도 등)가 성능 수준 자체를 결정함을 반영한다.

반면, 스코어링과 탐지 패밀리의 주효과는 각각 $\omega^2 = 0.027$, $0.025$로 작아, 스코어링이나 탐지 방법 자체보다 ``어떤 병원체에 어떤 방법을 적용하느냐''가 성능을 결정하는 지배적 요인임을 확인한다.

삼원 상호작용(score $\times$ family $\times$ pathogen)은 본 분석에서 보고하지 않았다.
이 항의 자유도는 $(9-1) \times (15-1) \times (9-1) = 896$으로, 전체 관측 수($n = 2{,}544$)에서 주효과 및 이원 상호작용의 자유도를 제외하면 잔차 자유도가 불충분하며, 불균형 설계(병원체별 유효 탐지 방법 수: 33--38개)로 인해 다수의 셀이 공백이므로 삼원 상호작용의 안정적 추정이 불가능하다.
이원 상호작용 중 family $\times$ pathogen($\omega^2 = 0.285$)이 이미 최대 분산 원천으로 확인되었으므로, 삼원 상호작용의 미보고가 핵심 결론에 영향을 미치지 않는다.

\subsection{LOPO 교차 검증}
\label{subsec:9pathogen_lopo}

8종 훈련 병원체에서 선정된 최적 조합이 제외된 1종 테스트 병원체에서도 최적인지를 LOPO(Leave-One-Pathogen-Out) 교차 검증으로 평가하였다.

\begin{itemize}
\item \textbf{조합 일치율}: 0/9 (0\%). 9회 반복 모두에서 훈련 최적 조합과 테스트 oracle 조합이 불일치하였다.
\item \textbf{평균 일반화 비율}: 4.9\%. 훈련 최적 조합의 테스트 MCC가 oracle MCC의 평균 4.9\%에 불과하였다.
\item \textbf{평균 성능 격차}: 0.277. 훈련 최적 조합과 oracle 간 MCC 차이의 평균이 0.277로, 최적 성능의 상당 부분을 손실하였다.
\item \textbf{최빈 훈련 최적 조합}: E\_only + SWAN이 8/9 반복에서 훈련 최적으로 선정되었으나, 이는 SWAN 탐지 방법이 Norovirus에서만 유효 결과(MCC = 0.311)를 산출하고 나머지 병원체에서는 MCC = 0.000을 기록하기 때문이다. 병원체별 유효 탐지 방법 수가 33--38개로 상이하여 SWAN 변이체가 모든 병원체에서 평가되지 않으며, 이로 인해 훈련 평균 MCC가 Norovirus 단일 병원체 성능에 의해 과대 추정된다.
\end{itemize}

이 결과는 ``다수 병원체에서 평균적으로 좋은 방법''을 선택하는 전략이 특정 병원체에 대한 예측에 실패함을 보여주며, 병원체별 적응적 방법 선택의 필요성을 강하게 지지한다.
SWAN 패밀리(9개 변이체)를 제외하고 LOPO를 수행하면\footnote{SWAN 제외 시 평가 수는 2,414개(원본 2,544개 대비 130개 감소)이며, 훈련 최적 조합이 주로 P*E + Wavelet($t$ = 1.5)으로 변경된다. 평균 일반화 비율은 4.9\%에서 11.6\%로 향상되고, 평균 성능 격차는 0.277에서 0.239로 감소하나, 조합 일치율 0/9는 유지되어 핵심 결론에 영향이 없다.}, 조합 일치율 0/9는 유지되어 ``병원체 간 최적 방법의 비이전성''이라는 핵심 결론은 변하지 않는다.

0/9의 조합 일치율과 4.9\%의 일반화 비율은 벤치마크의 한계가 아니라 \textbf{핵심 발견}이다. 이 결과는 ``단일 최적 방법''이라는 전제 자체가 성립하지 않음을 정량적으로 입증하며, 병원체 적응적 방법 선택(PAHD)의 필요성에 대한 가장 강력한 경험적 근거를 제공한다. $n=9$는 RNA 바이러스의 주요 대표 병원체를 포괄하며, 표본 크기의 제약에도 불구하고 병원체 간 이질성의 통계적 유의성이 ANOVA($\omega^2 = 0.285$), Friedman($p = 0.0015$), LOPO(0/9 일치) 세 가지 독립적 검정에서 모두 확인되었다.

\subsection{Friedman 검정}
\label{subsec:9pathogen_friedman}

9개 병원체를 블록으로 사용하여, 상위 20개 스코어링--탐지 조합 간 성능 차이에 대한 Friedman 검정을 수행하였다.
$k = 20$은 전체 유효 조합 수(병원체별 247--305개)의 상위 약 7\%에 해당하며, 경쟁력 있는 조합군을 대표하면서도 블록 수($b = 9$)에 비해 과도하지 않은 처리 수를 유지한다\footnote{$k$ 선택의 강건성을 확인하기 위해, $k = 10$, $k = 30$, $k = 50$에서도 Friedman 검정을 수행하였다. $k = 10$: $\chi^2 = 22.1$, $p = 0.009$; $k = 30$: $\chi^2 = 58.7$, $p = 0.001$; $k = 50$: $\chi^2 = 87.3$, $p < 0.001$. 모든 $k$에서 $p < 0.01$로 유의하여, 상위 조합 간 성능 차이의 유의성은 $k$ 선택에 강건하다.}.

\begin{equation}
\chi^2 = 42.44, \quad p = 0.0015
\end{equation}

$p = 0.0015 < 0.01$로 통계적으로 유의하여, 상위 20개 조합 간에 유의한 성능 차이가 존재함을 확인하였다.
이는 ``상위 조합들이 모두 동등하므로 아무 조합이나 사용해도 된다''는 귀무 가설을 기각하며, 조합 선택이 성능에 실질적 영향을 미침을 입증한다.

\subsection{병원체 수준 BCa 부트스트랩 신뢰 구간}
\label{subsec:9pathogen_bootstrap}

전체 병원체에 걸친 최적 단일 조합을 BCa 부트스트랩으로 평가하였다.
선정 기준은 95\% BCa 신뢰 구간의 하한이 0을 초과하는 조합 중 평균 MCC가 가장 높은 조합이다.
병원체별 평균 MCC만으로 순위를 매기면 특정 병원체에서만 높은 성능을 보이는 조합(예: E\_only + SWAN은 Norovirus에서 MCC = 0.311이나 나머지 8개 병원체에서 MCC = 0.000)이 상위에 위치할 수 있으므로, 신뢰 구간이 0을 제외하는 조건을 통해 전체 병원체에 걸친 일관된 성능을 보장한다.

\begin{itemize}
\item \textbf{최적 단일 조합}: P*E + Wavelet ($t$ = 1.5)
\item \textbf{평균 MCC}: 0.124
\item \textbf{95\% BCa 신뢰 구간}: [0.047, 0.202]
\end{itemize}

신뢰 구간의 폭(0.155)이 넓으며, 이는 9개 병원체에 걸친 성능의 높은 변동성을 반영한다.
최적 단일 조합의 평균 MCC(0.124)가 병원체별 oracle MCC의 평균(0.298)의 42\%에 불과하여, 단일 고정 조합으로는 전체 병원체의 절반 이상에서 최적 성능에 크게 미달함을 정량적으로 확인한다.

이 MCC 값은 낮아 보이지만, 다음 세 가지 맥락에서 해석해야 한다. 첫째, 이는 단일 고정 조합의 9개 병원체 평균이며, 병원체별 oracle MCC(평균 0.298)의 42\%에 해당한다. 병원체별 최적 조합을 사용하면 MCC가 0.088--0.549 범위로 실질적으로 향상된다. 둘째, 무작위 기준선(MCC 평균=0.001, 표준편차=0.003)과 비교하면, 관찰된 MCC=0.124는 무작위 기대를 약 41 표준편차 초과한다. 이 $z$-점수 해석은 무작위 MCC 분포의 근사적 정규성을 가정하며, 실제 분포가 정규성에서 벗어나더라도 Chebyshev 부등식에 의해 $P(|X - \mu| \geq 41\sigma) \leq 1/41^2 \approx 0.06\%$이므로 통계적 유의성은 유지된다. 셋째, 핫스팟 탐지는 본질적으로 불균형 이진 분류(양성 비율 약 1--16\%)이며, 이러한 극단적 불균형에서 MCC $>$ 0.1은 무의미하지 않다.

\subsection{하위표본 강건성}
\label{subsec:9pathogen_subsampling}

MSA 서열의 25\%, 50\%, 75\%, 100\% 비율에서 벤치마크를 반복 수행한 결과, 전체 병원체 평균 MCC는 0.054--0.058 범위에서 안정적이었다.
표본 크기 감소에 따른 MCC 변동이 $\pm$0.002 이내로, 벤치마크 결과가 서열 표본 크기에 강건함을 확인하였다.
이는 본 벤치마크의 결론이 특정 표본 크기의 인공물이 아닌 병원체 변이 구조의 본질적 특성을 반영함을 뒷받침한다.

\subsection{핵심 발견 요약: 병원체 적응적 방법 선택의 필요성}
\label{subsec:pahd_motivation_summary}

9-병원체 벤치마크의 핵심 결과를 종합하면, 병원체 적응적 방법 선택의 필요성이 네 가지 독립적 증거에 의해 입증된다.

\begin{enumerate}
\item \textbf{상호작용 지배}: 탐지 패밀리 $\times$ 병원체 상호작용($\omega^2 = 0.285$)이 모든 주효과보다 크다. 최적 방법은 병원체에 의존한다.
\item \textbf{고유 최적 조합}: 9/9 병원체에서 9개의 고유한 최적 조합이 관찰되었다. 두 병원체가 동일한 최적 조합을 공유하는 경우는 없다.
\item \textbf{일반화 실패}: LOPO 교차 검증에서 0/9 조합 일치, 평균 일반화 비율 4.9\%. 훈련 병원체에서 선정된 최적 조합은 새로운 병원체에 일반화되지 않는다.
\item \textbf{유의한 조합 차이}: Friedman $\chi^2 = 42.44$, $p = 0.0015$. 조합 선택은 성능에 유의한 영향을 미친다.
\end{enumerate}

이러한 결과는 ``병원체마다 최적 방법이 다르다''는 경험적 사실을 통계적으로 확립하며, 병원체의 변이 특성을 사전에 프로파일링하여 최적 조합을 적응적으로 선택하는 PAHD 프레임워크의 동기를 제공한다.

\subsection{PAHD 프로토타입 검증}
\label{subsec:pahd_prototype}

앞선 LOPO 분석에서 naive 전략(훈련 병원체의 평균 최적 조합을 그대로 적용)이 0/9 일치, 평균 일반화 비율 4.9\%로 완전히 실패하였다.
이에 PAHD의 핵심 가설---``변이 프로파일이 유사한 병원체는 유사한 최적 탐지 조합을 가진다''---을 검증하기 위한 프로토타입 실험을 수행하였다.

\subsubsection{병원체 프로파일 구성}

두 가지 버전의 병원체 프로파일을 평가하였다.
\textbf{v1 프로파일}은 게놈 및 ground truth 특성에서 추출한 8개 특성 벡터이다(oracle 성능 정보 미포함): 정렬 길이, 고유 서열 수, Layer~A(양성) ground truth 크기, Layer~B(음성) ground truth 크기, 다양성 비율(고유 서열 수/전체 서열 수), 적응적 밀도(Layer~A 크기/정렬 길이), 구속 밀도(Layer~B 크기/정렬 길이), ground truth 비율(Layer~A / (Layer~A + Layer~B)).

\textbf{v2 프로파일}은 MSA 계산 점수 분포에서 85개 특성을 추출한다: 스코어링 공식별 통계량(평균, 표준편차, 왜도, 첨도, 영 비율, Gini 계수), 교차 점수 상관 구조, 공간적 자기상관(Moran's~I) 등이며, 모두 oracle 성능 정보 없이 계산된다.
v2 프로파일은 v1에서 관찰된 병리적으로 높은 유사도를 보정하여(예: MERS--SARS-CoV-2 코사인 유사도가 0.834에서 0.170으로 감소), 메타데이터 수준 특성이 아닌 실제 점수 분포 특성을 포착함으로써 생물학적으로 더 타당한 병원체 그룹핑을 산출한다.

두 프로파일 모두 z-표준화 후 코사인 유사도로 비교하였다.

\subsubsection{프로파일 기반 LOPO 예측}

세 가지 PAHD 전략을 비교하였다: 1-NN(가장 유사한 병원체의 최적 조합), Weighted(유사도 가중 MCC 투표), Family-first(2단계: 최적 탐지 패밀리 예측 후 해당 패밀리 내 최적 조합 선택).
표~\ref{tab:pahd_prototype}은 최우수 전략인 PAHD Family-first의 v1 및 v2 프로파일 결과를 naive LOPO와 비교한다.

\begin{table}[htbp]
\centering
\caption{Naive LOPO 대 PAHD Family-first 프로토타입의 병원체별 테스트 MCC 비교 (v1: 8개 특성, v2: 85개 MSA 유래 특성). Oracle은 각 병원체의 최적 조합 MCC이다. 프로파일 특성은 순수 게놈 특성만 사용(oracle 성능 정보 미포함).}
\label{tab:pahd_prototype}
\adjustbox{max width=\textwidth}{
\begin{tabular}{lrrrr}
\toprule
\textbf{병원체} & \textbf{Oracle} & \textbf{Naive} & \textbf{PAHD Family v1} & \textbf{PAHD Family v2} \\
\midrule
Dengue      & 0.256 & 0.000 & $-$0.012 & $-$0.027 \\
H3N2        & 0.549 & 0.000 & 0.207 & 0.011 \\
HCV         & 0.443 & 0.000 & $-$0.291 & 0.141 \\
HIV-1       & 0.182 & 0.000 & 0.014 & 0.014 \\
Influenza B & 0.361 & 0.000 & $-$0.190 & $-$0.027 \\
MERS        & 0.088 & 0.000 & $-$0.012 & $-$0.025 \\
Norovirus   & 0.438 & 0.192 & 0.217 & 0.118 \\
RSV         & 0.249 & 0.000 & 0.069 & $-$0.054 \\
SARS-CoV-2  & 0.118 & 0.000 & 0.045 & 0.000 \\
\midrule
평균         & 0.298 & 0.021 & 0.005 & 0.017 \\
양성 MCC 수  & 9/9   & 1/9   & \textbf{5/9} & 4/9 \\
\bottomrule
\end{tabular}
}
\end{table}

비교를 위해 세 가지 PAHD 전략을 v1 및 v2 프로파일 모두에서 평가하였다. v1 프로파일에서: 1-NN은 3/9 양성 MCC(평균 $= -0.027$), Weighted는 1/9(평균 $= 0.024$), Family-first는 최고의 양성 MCC 수(5/9, 평균 $= 0.005$)를 달성하였다. v2 프로파일(85개 MSA 유래 특성)에서: 1-NN은 3/9 양성 MCC(평균 $= -0.016$), Weighted는 2/9(평균 $= 0.024$), Family-first는 4/9(평균 $= 0.017$)를 달성하였다.

v2 프로파일의 주된 개선은 Family-first 전략의 평균 MCC가 실질적으로 향상된 것이다($0.005 \rightarrow 0.017$). 이는 병리적 음의 전이를 제거한 결과이다: HCV가 $-0.291$에서 $+0.141$로, Influenza~B가 $-0.190$에서 $-0.027$로 개선되었다. 이러한 보정은 v2 프로파일이 생물학적으로 상이한 병원체를 혼동하지 않는 더 판별적인 유사도 점수를 산출하기 때문이다.

\subsubsection{결과 해석}

PAHD Family-first 전략의 주요 이점은 \textbf{양성 전이의 폭}이다: 양성 MCC를 산출하는 병원체 수를 1/9(naive)에서 5/9(v1) 또는 4/9(v2)로 증가시켜, 프로파일 기반 방법 선택이 탐지 능력을 단일 유리한 경우에 집중시키기보다 병원체 전반에 분산시킴을 보여준다.
\textbf{중위수 MCC}가 이 패턴을 확인한다: PAHD v1은 중위수 0.014를 달성한 반면 naive의 중위수는 0.000이며, 이는 9개 중 8개 병원체가 naive에서 MCC = 0을 받기 때문이다.
naive 평균 MCC(0.021)는 단일 Norovirus 이상치(MCC $= 0.192$)에 의해 부풀려지며; 이 이상치를 제외하면 naive 평균은 $-0.003$으로 하락한다.

v2 프로파일은 다른 차원의 개선을 제공한다: \textbf{병리적 음의 전이 감소}. v1이 더 많은 양성 MCC 사례를 달성하지만(5/9 vs.\ 4/9), 평균 MCC($0.005$)가 큰 음의 이상치(HCV $= -0.291$, Influenza~B $= -0.190$)에 의해 끌려 내려간다. v2 프로파일은 이러한 극단적 실패를 제거하여(HCV $= +0.141$, Influenza~B $= -0.027$), 양성 사례 수는 적지만 더 높은 평균 MCC $0.017$을 산출한다.
이 트레이드오프는 v2 프로파일의 개선된 유사도 구조를 반영한다: 인위적으로 높은 병원체 간 유사도를 보정함으로써, v2는 비유사한 병원체로부터의 탐지 패밀리 전이를 회피하여 치명적 실패를 줄이되 일부 보수적 예측을 감수한다.

탐지 패밀리 수준의 top-$k$ 적중 기준으로 평가할 때, PAHD Family-first 전략은 패밀리 top-3 적중률 5/9를 달성하였다. 즉, 9개 병원체 중 5개에서 예측된 탐지 패밀리가 해당 병원체의 상위 3개 고성능 패밀리에 포함되었다. 정확한 조합 수준에서는 top-10 적중률이 1/9(H3N2, 369개 중 8위)로, 정확한 조합 예측은 9개 병원체만으로는 어렵지만 예측된 탐지 패밀리가 oracle의 경쟁적 근방에 위치하는 경우가 빈번함을 확인하였다.

9개 훈련 인스턴스만으로 351개 정확한 조합을 예측하는 비적정 문제를 추가로 분석하기 위해, v2 프로파일 유사도 행렬과 1-NN LOPO를 사용하여 더 거친 예측 단위에서도 PAHD를 평가하였다.
\textbf{탐지 패밀리 수준}(15개 패밀리)에서 1-NN 전략은 패밀리 일치율 0/9를 기록하였다---각 병원체의 oracle 최적 패밀리가 고유하며, 1-NN 이웃은 동일한 최적 패밀리를 공유하지 않는다.
그러나 예측된 패밀리 \emph{내}에서 최적 조합을 사용할 때의 평균 복합 MCC는 0.144(중위수 0.142)로, oracle 성능의 49.2\%를 회복하였다.
\textbf{스코어링 공식 수준}(9개 공식)에서도 일치율은 0/9이었으나, 평균 MCC는 0.172(중위수 0.138)로 oracle 성능의 54.5\%를 회복하였다.
두 거친 단위 예측 모두 정확한 조합 예측(평균 MCC $= -0.016$, oracle의 $-9.8$\%)과 naive 전역 최적 기준선(평균 MCC $= 0.034$, oracle의 8.2\%)을 실질적으로 능가한다.
핵심 통찰은, \emph{라벨} 예측(어떤 패밀리 또는 스코어링)이 불충분한 $n$으로 인해 실패하더라도, 예측된 범주를 적용하여 달성되는 \emph{성능}은 경쟁력을 유지한다는 것이다: 검색 공간을 단일 예측 패밀리 또는 스코어링 공식으로 제한하는 것만으로도 oracle MCC의 약 절반을 포착하며, naive 접근법에서는 거의 0에 가깝다.

한계는 두 가지이다.
첫째, 9개 병원체의 knowledge base는 프로파일 공간에서의 일반화에 불충분하다.
v1에서 가장 큰 음의 MCC를 보인 두 병원체---HCV($-$0.291)와 Influenza~B($-$0.190)---는 프로파일 공간에서 유사한 이웃이 없기 때문이다.
v2 프로파일은 점수 분포 수준 특성을 포착하여 이러한 실패를 실질적으로 완화하지만(HCV $-0.291 \rightarrow +0.141$), 모든 음의 전이를 제거하지는 못한다.
둘째, 정확한 조합 일치율은 v1과 v2 프로파일 모두에서 0/9으로, 85개 MSA 유래 특성을 사용하더라도 9개 병원체만으로는 단일 최적 조합 예측이 지나치게 야심적임을 확인하며, 상위 $k$개 후보 풀에서의 적응적 필터링이 더 현실적인 전략이다.
v2 프로파일의 개선에도 불구하고, $n = 9$ 병원체의 근본적 한계가 지속된다: LOPO 일치율은 0/9로 유지되어, 신뢰할 수 있는 조합 예측을 위해서는 knowledge base를 9개 병원체 이상으로 확대해야 함을 확인한다.

이러한 결과는 9개 병원체와 oracle 정보 없이도 게놈 프로파일 유사도가 탐지 패밀리 선호도를 부분적으로 예측할 수 있음(v1: 5/9 양성 MCC; v2: 평균 MCC $0.005 \rightarrow 0.017$, 음의 전이 감소)을 보여준다.
Knowledge base의 규모 확대(20개 이상), 가중 앙상블 전략 도입, 메타학습 회귀 모델~\cite{rice1976algorithm,smith2021algorithm}로의 전환이 실용적 PAHD 구현에 필수적이다.

\subsection{ANOVA 진단 검정}
\label{subsec:anova_diagnostics}

Shapiro-Wilk 검정에서 잔차의 정규성이 기각되었으나($W = 0.971$, $p = 3.29 \times 10^{-22}$), 왜도 = 0.12, 첨도 = 2.42로 경미한 비정규성이며, Levene 검정에서도 모든 요인의 등분산성이 기각되었다.
그러나 대표본($n = 2{,}544$)에서 ANOVA는 이러한 위반에 강건하며, Kruskal-Wallis 비모수 검정(병원체: $H = 569.1$, $p < 10^{-117}$; 탐지 패밀리: $H = 86.1$, $p < 10^{-12}$; 스코어링: $H = 78.1$, $p < 10^{-13}$)에서 모든 효과가 동일하게 확인되어 ANOVA 결론의 신뢰성이 유지된다.

\subsection{Precision@$k$ 분석}
\label{subsec:precision_at_k}

병원체별 최적 조합의 탐지 정밀도와 재현율을 분석하였다. 표~\ref{tab:precision_at_k}은 각 병원체의 최적 조합(top-1) 기준 정밀도, 재현율, 탐지 크기, 최고 MCC를 제시한다.

\begin{table}[htbp]
\centering
\caption{병원체별 최적 조합(top-1)의 정밀도, 재현율, 탐지 크기 및 최고 MCC.}
\label{tab:precision_at_k}
\begin{tabular}{lrrrr}
\toprule
병원체 & Precision & Recall & $n_{\text{det}}$ & Best MCC \\
\midrule
Dengue       & 0.157 & 0.667 & 102 & 0.256 \\
H3N2         & 0.333 & 1.000 &  75 & 0.549 \\
HCV          & 0.622 & 0.426 &  37 & 0.443 \\
HIV-1        & 0.089 & 0.913 & 235 & 0.182 \\
Influenza B  & 0.189 & 0.824 &  74 & 0.361 \\
MERS         & 0.018 & 0.778 & 390 & 0.088 \\
Norovirus    & 0.290 & 0.733 &  38 & 0.438 \\
RSV          & 0.667 & 0.100 &   3 & 0.249 \\
SARS-CoV-2   & 0.090 & 0.240 &  67 & 0.118 \\
\midrule
평균          & 0.273 & 0.631 &  -- & 0.298 \\
\bottomrule
\end{tabular}
\end{table}

전체 병원체의 top-1 평균 정밀도는 0.273, 평균 재현율은 0.631이다.
H3N2는 완전한 재현율(1.0)과 양호한 정밀도(0.333)를 보여 탐지가 가장 용이한 병원체이며, RSV는 가장 높은 정밀도(0.667)를 기록하였으나 재현율이 0.10으로 극단적인 트레이드오프를 나타낸다.
MERS는 ground truth 크기(9개 위치) 대비 탐지 크기(390개)가 과도하게 커 정밀도가 0.018로 가장 낮다.
이러한 병원체 간 정밀도--재현율 프로파일의 이질성은 단일 임계값 기반 탐지의 한계를 추가로 확인한다.

\subsection{DMS-진화 상관 분석}
\label{subsec:dms_evolution_correlation}

DMS 적합도 데이터가 존재하는 4개 병원체에 대해, 위치별 DMS fitness 효과와 최적 스코어링 출력값 간 Spearman 상관을 분석하였다(표~\ref{tab:dms_evolution_correlation}).

\begin{table}[htbp]
\centering
\caption{DMS 적합도 점수와 진화적 탐지 점수 간 Spearman 상관.}
\label{tab:dms_evolution_correlation}
\begin{tabular}{lrr}
\toprule
병원체 & $\rho$ & $p$ \\
\midrule
H3N2       &  0.191 & $1.5 \times 10^{-3}$ \\
HIV-1      &  0.412 & $< 10^{-12}$ \\
RSV        & $-$0.363 & $< 10^{-9}$ \\
SARS-CoV-2 & $-$0.370 & $< 10^{-10}$ \\
\bottomrule
\end{tabular}
\end{table}

상관 방향이 병원체에 따라 상이하다---H3N2와 HIV-1은 양의 상관, RSV와 SARS-CoV-2는 음의 상관---이는 DMS 적합도와 진화적 적응이 변이 영향의 근본적으로 다른 측면을 측정함을 확인한다.
이 결과는 3-계층 ground truth 프레임워크가 DMS와 진화적 증거를 독립적이고 상호 보완적인 계층으로 처리하는 설계 원칙을 검증한다.

\subsection{DMS Ground Truth 역치 민감도 분석}
\label{subsec:dms_threshold_sensitivity}

Layer C (DMS)의 ``상위 20백분위'' 역치가 자의적이므로, 10\%, 20\%(기본값), 30\% 역치에서 DMS ground truth를 재구성하고 4종 병원체에서 벤치마크를 반복하였다(표~\ref{tab:dms_threshold_sensitivity}).
4/4 병원체에서 최적 스코어링--탐지 조합이 역치에 따라 변경되어, DMS 기반 평가에서 역치 선택이 방법 순위에 영향을 미치며 투명한 보고가 필요함을 확인하였다.

\begin{table}[htbp]
\centering
\caption{DMS ground truth 역치에 따른 최적 스코어링--탐지 조합과 DMS F1 (4종 병원체)}
\label{tab:dms_threshold_sensitivity}
\small
\begin{tabular}{llllr}
\toprule
\textbf{역치} & \textbf{병원체} & \textbf{$n_{\text{GT}}$} & \textbf{최적 조합} & \textbf{DMS F1} \\
\midrule
\multirow{4}{*}{10\%}
& H3N2 & 52 & rank(P*E) + Wavelet($t$=1.0) & 0.443 \\
& HIV-1 & 54 & E*rare + KDE($p$=80) & 0.347 \\
& RSV & 51 & rank(P*E) + KDE($p$=75) & 0.275 \\
& SARS-CoV-2 & 126 & rank(P*E) + KDE($p$=80) & 0.249 \\
\midrule
\multirow{4}{*}{20\% (기본)}
& H3N2 & 102 & rank(P*E) + AdaptWin($s$=90) & 0.434 \\
& HIV-1 & 100 & P*E*rare + KDE($p$=80) & 0.400 \\
& RSV & 101 & P*minority\_E + MutClust-Orig & 0.320 \\
& SARS-CoV-2 & 252 & E*rare + ScoreDBSCAN($\epsilon$=15) & 0.353 \\
\midrule
\multirow{4}{*}{30\%}
& H3N2 & 156 & rank(P*E) + AdaptiveKnee & 0.499 \\
& HIV-1 & 144 & minority\_E + AdaptWin($s$=90) & 0.498 \\
& RSV & 151 & P*minority\_E + MutClust-Orig & 0.438 \\
& SARS-CoV-2 & 377 & E*rare + ScoreDBSCAN($\epsilon$=15) & 0.483 \\
\bottomrule
\end{tabular}
\begin{minipage}{\textwidth}
\smallskip
\footnotesize $n_{\text{GT}}$: 해당 역치에서의 DMS ground truth 위치 수. 평균 최고 DMS F1은 10\%: 0.328, 20\%: 0.377, 30\%: 0.479로 역치가 느슨할수록 증가한다. 4/4 병원체에서 최적 스코어링--탐지 조합이 역치에 따라 변경되었다.
\end{minipage}
\end{table}


%% ============================================================
\section{삽입-결실 변이(Indel) 배제의 영향 분석}
\label{sec:indel_analysis}

H-score는 참조 서열 대비 각 고정 좌표의 점 변이 빈도를 측정하는 고정 좌표 기반 지표이므로, 삽입(insertion) 및 결실(deletion) 변이는 구조적으로 포착할 수 없다.
GISAID 2024--2025 Spike 서열 9,957개를 분석한 결과, 99.8\%가 최소 1개의 기능적 indel(예: del69-70, Y144del, del156-157)을 보유하며, 모든 indel 위치의 H-score는 정확히 0.000으로 H-score 기반 탐지에서 완전히 배제된다.
약 35개의 주요 Spike 기능적 변이 중 최소 5개(약 14\%)가 indel이다.

탐지 영향을 정량화하기 위해, 1,251 nt ground truth에 5개 indel 위치(27 nt)를 추가한 확장 ground truth를 구성하였다.
원래 H-score와 augmented H-score(indel 위치에 기능적 신호 주입)를 비교한 합성 실험에서 MutClust-Hybrid의 F1 영향은 $\Delta$F1 = $-$0.002로 미미하였는데, 이는 27개 추가 위치가 확장 ground truth의 2.1\%에 불과하기 때문이다.
F1에 대한 정량적 영향은 작지만, 생물학적으로 중요한 구조적 변이가 스코어링 함수에서 체계적으로 보이지 않는다는 점에서 고정 좌표 프레임워크의 구조적 한계를 나타낸다.
